\documentclass[12pt]{article}
\usepackage[utf8]{inputenc}
\usepackage[T1]{fontenc}
\usepackage{hyperref}
\title{S\'enoufo, Tagwana Language Resources for K-12 Curriculum Development}
\author{Compiled by Claude Sonnet 5 (claude-sonnet-5) for Andr\'as Kornai}
\date{}
\begin{document}
\maketitle

\section{Introduction}
Tagwana (ISO 639-3 tgw) is a Senufo language of northern C\^ote d'Ivoire, in the Tagwana--Djimini branch. Global Recordings lists the language as ``Tagouana'' with a population of 138,100, and gives local varieties such as Fourgoula, Katiola and Niakaramadougou. The search found no school materials. Tagwana is sometimes treated as a separate language and sometimes as a dialect of Senufo, and I found no source that settles this.

\section{Primary Educational Resources}
None were found in search results.

\section{Linguistic and Reference Materials}
\begin{itemize}
\item Tagwana--Djimini languages, Wikipedia: \url{https://en.wikipedia.org/wiki/Tagwana%E2%80%93Djimini_languages}
\item A 2021 article in \emph{Studies in African Linguistics} on Tagbana tone, a closely related Senufo language. Listed by DOAJ: \url{https://doaj.org/article/6bff25ec8de7409f9dfc1354a5068995}
\end{itemize}

\section{Religious and Community Resources}
Joshua Project profile of the Tagwana Senoufo: \url{https://joshuaproject.net/people_groups/15181/IV}

\section{Digital and Audio Resources}
Global Recordings Network lists the language under the Fourgoula variety; the page presumably holds audio recordings, though the search result showed no details: \url{https://globalrecordings.net/en/language/16534}

\section{Curriculum Development Context}
The language has many speakers but no school materials turned up. A curriculum would need a standardized orthography and a decision on which variety to use. It could build on materials from related Senufo languages, if such materials exist.
\end{document}